%! Simplifying Radicals; Adding \& Subtracting — Unit 5, Lesson 5.2 — Homework (Answer Key)
\documentclass[10pt]{article}
\usepackage{algebra2-article}
\usepackage{algebra2-key}   % pulls in -boxes; do NOT also load -boxes
\newcommand{\TallMath}[1]{$\displaystyle #1\rule[-1.4em]{0pt}{3.2em}$}
% In-formula answer slot (\ans is text-mode and must never appear inside $...$).
\newcommand{\mans}[1]{{\color{keyred}\mathbf{#1}}}

\begin{document}

\pageheader{Unit 5, Lesson 5.2}{Homework --- Answer Key}
\namedateperiod

\vspace{0.10in}

\begin{notesbox}{Practice}
Assume all variables represent \textbf{positive} real numbers. No calculator except where a problem
asks for a decimal.

\begin{enumerate}[leftmargin=1.9em, itemsep=11pt]
  \item \textbf{Simplify each numeric radical.} Watch the index --- it tells you the group size.
    \begin{center}
    \renewcommand{\arraystretch}{1.9}
    \begin{tabularx}{\linewidth}{@{}X X X X@{}}
    (a) \; $\sqrt{28}=$ \ans{$2\sqrt7$} & (b) \; $\sqrt{96}=$ \ans{$4\sqrt6$} &
    (c) \; $\sqrt{147}=$ \ans{$7\sqrt3$} & (d) \; $\sqrt{242}=$ \ans{$11\sqrt2$} \\
    (e) \; $\sqrt[3]{40}=$ \ans{$2\sqrt[3]{5}$} & (f) \; $\sqrt[3]{-135}=$ \ans{$-3\sqrt[3]{5}$} &
    (g) \; $\sqrt[4]{80}=$ \ans{$2\sqrt[4]{5}$} & (h) \; $\sqrt[5]{64}=$ \ans{$2\sqrt[5]{2}$} \\
    \end{tabularx}
    \end{center}
    \emph{Work:} (a) $4\cdot7$. \; (b) $16\cdot6$. \; (c) $49\cdot3$. \; (d) $121\cdot2$. \;
    (e) $8\cdot5$. \; (f) $-27\cdot5$ --- an odd index accepts a negative radicand (Lesson 5.1). \;
    (g) $16\cdot5$. \; (h) $32\cdot2$.

  \item \textbf{Simplify each algebraic radical.} Divide the exponent by the index: quotient out,
        remainder in.
    \begin{center}
    \renewcommand{\arraystretch}{2.0}
    \begin{tabularx}{\linewidth}{@{}X X X@{}}
    (a) \; $\sqrt{x^{11}}=$ \ans{$x^{5}\sqrt{x}$} & (b) \; $\sqrt{20a^{3}}=$ \ans{$2a\sqrt{5a}$} &
    (c) \; $\sqrt{72m^{4}n^{5}}=$ \ans{$6m^{2}n^{2}\sqrt{2n}$} \\
    (d) \; $\sqrt[3]{24p^{7}}=$ \ans{$2p^{2}\sqrt[3]{3p}$} &
    (e) \; $\sqrt[3]{-16c^{9}}=$ \ans{$-2c^{3}\sqrt[3]{2}$} &
    (f) \; $\sqrt{\dfrac{18y^{6}}{25}}=$ \ans{$\dfrac{3y^{3}\sqrt2}{5}$} \\
    \end{tabularx}
    \end{center}
    \emph{Work:} (a) $11\div2$ is $5$ r$1$. \; (b) $20=4\cdot5$, $a^{3}=a^{2}\cdot a$. \;
    (c) $72=36\cdot2$, $m^{4}$ comes out whole, $n^{5}=n^{4}\cdot n$. \;
    (d) $24=8\cdot3$, $7\div3$ is $2$ r$1$. \; (e) $-16=-8\cdot2$, $9\div3$ is $3$ r$0$. \;
    (f) $\sqrt{18}=3\sqrt2$ over $\sqrt{25}=5$, and $y^{6}$ comes out whole.

  \item \textbf{Add or subtract.} Simplify every term first. If a sum cannot be combined, say so ---
        but only after you have simplified.
    \begin{center}
    \renewcommand{\arraystretch}{2.0}
    \begin{tabularx}{\linewidth}{@{}X X@{}}
    (a) \; $6\sqrt{11}-2\sqrt{11}=$ \ans{$4\sqrt{11}$} &
    (b) \; $\sqrt{18}+\sqrt{32}=$ \ans{$7\sqrt2$} \\
    (c) \; $\sqrt{50}-\sqrt{8}+\sqrt{2}=$ \ans{$4\sqrt2$} &
    (d) \; $2\sqrt{63}+4\sqrt{28}=$ \ans{$14\sqrt7$} \\
    (e) \; $\sqrt{20}+\sqrt{12}=$ \ans{$2\sqrt5+2\sqrt3$ --- cannot combine} &
    (f) \; $\sqrt[3]{54}-\sqrt[3]{2}=$ \ans{$2\sqrt[3]{2}$} \\
    (g) \; $\sqrt{45x}-\sqrt{20x}=$ \ans{$\sqrt{5x}$} &
    (h) \; $\sqrt{8}+\sqrt[3]{8}=$ \ans{$2\sqrt2+2$ --- cannot combine} \\
    \end{tabularx}
    \end{center}
    \emph{Work:} (b) $3\sqrt2+4\sqrt2$. \; (c) $5\sqrt2-2\sqrt2+\sqrt2$. \;
    (d) $2(3\sqrt7)+4(2\sqrt7)=6\sqrt7+8\sqrt7$. \; (e) both simplify, but to \emph{different}
    radicands. \; (f) $3\sqrt[3]{2}-\sqrt[3]{2}$. \; (g) $3\sqrt{5x}-2\sqrt{5x}$. \;
    (h) same radicand $8$, but the \emph{indices} differ --- and $\sqrt[3]{8}$ is just the number
    $2$, so the answer is already as simple as it gets.

  \item \textbf{Use the quotient property.}
    \begin{center}
    \renewcommand{\arraystretch}{2.0}
    \begin{tabularx}{\linewidth}{@{}X X X@{}}
    (a) \; $\sqrt{\dfrac{121}{9}}=$ \ans{$\dfrac{11}{3}$} &
    (b) \; $\dfrac{\sqrt{54}}{\sqrt{6}}=$ \ans{$3$} &
    (c) \; $\sqrt{\dfrac{5}{9}}=$ \ans{$\dfrac{\sqrt5}{3}$} \\
    \end{tabularx}
    \end{center}
    \emph{Note on (b):} run the property backwards --- $\sqrt{54/6}=\sqrt{9}=3$ is far easier than
    simplifying each radical separately.

  \item \textbf{Two questions about what ``simplest'' means.}
    \begin{enumerate}[label=(\alph*), leftmargin=1.8em, itemsep=6pt]
      \item A classmate claims $\sqrt{36+64}=\sqrt{36}+\sqrt{64}$. Evaluate \emph{both} sides to
            show they are not equal, then write the rule that \emph{does} work and check it on the
            same two numbers. \ansline{Left: $\sqrt{36+64}=\sqrt{100}=10$. Right:
            $\sqrt{36}+\sqrt{64}=6+8=14$. Since $10\ne14$, $\sqrt{a+b}\ne\sqrt a+\sqrt b$.\\[3pt]
            The rule that works is the \textbf{product} property, $\sqrt{ab}=\sqrt a\cdot\sqrt b$:
            $\sqrt{36\cdot64}=\sqrt{2304}=48$, and $\sqrt{36}\cdot\sqrt{64}=6\cdot8=48$. A radical
            splits over multiplication, never over addition.}
      \item Explain why $\sqrt[3]{16}$ is \emph{not} in simplest form but $\sqrt[3]{9}$ \emph{is} ---
            even though $9$ is a perfect square. \ansline{Under a \emph{cube} root the only powers
            that come out are perfect \emph{cubes}. $16=8\cdot2$ contains the perfect cube $8$, so
            $\sqrt[3]{16}=2\sqrt[3]{2}$ --- not finished. $9=3\cdot3$ contains no perfect cube (you
            would need \emph{three} equal factors, and there are only two), so it is already in
            simplest form. Being a perfect square is irrelevant here: the index decides what counts.}
    \end{enumerate}

  \item \textbf{Two roads, one answer (connecting to Lesson 5.1).} Simplify $\sqrt{x^{9}}$ twice.
    \begin{enumerate}[label=(\alph*), leftmargin=1.8em, itemsep=5pt]
      \item By splitting the radicand: $\sqrt{x^{9}}=\sqrt{x^{\mans{8}}\cdot x}=$ \ans{$x^{4}\sqrt{x}$}
      \item With a rational exponent: $\sqrt{x^{9}}=x^{\mans{9/2}}=x^{\mans{4}}\cdot
            x^{\mans{1/2}}=$ \ans{$x^{4}\sqrt{x}$}
      \item Why do these two methods always agree? Use the words \emph{mixed number} in your
            answer. \ansline{Because dividing the exponent by the index \emph{is} rewriting
            $\frac{9}{2}$ as the \textbf{mixed number} $4\frac12$. The whole-number part, $4$,
            becomes the power that comes outside; the leftover $\frac12$ is exactly one square
            root's worth and stays inside. Splitting $x^{9}$ into $x^{8}\cdot x$ is the same
            division written without fractions.}
    \end{enumerate}
\end{enumerate}
\end{notesbox}

\begin{extensionbox}
\textbf{The pendulum: quadruple the length to double the swing.} A pendulum of length $L$ feet takes
$T$ seconds to complete one full swing, where
\[
  T(L) = 2\pi\sqrt{\frac{L}{32}}\ \text{ seconds.}
\]
Every length below was chosen so that the fraction under the radical reduces to a perfect square.
\begin{center}
\renewcommand{\arraystretch}{1.3}
\begin{tabular}{|c|c|c|c|c|}
\hline
$L$ (feet)               & $2$ & $8$ & $18$ & $32$ \\ \hline
$\sqrt{L/32}$ (reduced)  & $\frac14$ & \ans{$\frac12$} & \ans{$\frac34$} & \ans{$1$} \\ \hline
$T(L)$ (exact)           & $\frac{\pi}{2}$ & \ans{$\pi$} & \ans{$\frac{3\pi}{2}$} & \ans{$2\pi$} \\ \hline
\end{tabular}
\end{center}
\begin{enumerate}[label=(\alph*), leftmargin=1.8em, itemsep=5pt]
  \item Complete the table \emph{exactly}, leaving $\pi$ in your answers. (Reduce the fraction
        \emph{before} taking the root --- simplest-form condition 2 says no fractions under a
        radical.) \ansline{$\frac{8}{32}=\frac14$ and $\sqrt{\frac14}=\frac12$;
        $\frac{18}{32}=\frac{9}{16}$ and $\sqrt{\frac{9}{16}}=\frac34$; $\frac{32}{32}=1$. Multiply
        each by $2\pi$.}
  \item Give each period as a decimal, rounded to the nearest hundredth of a second.
        \ansline{$1.57$ s, \; $3.14$ s, \; $4.71$ s, \; $6.28$ s.}
  \item From $L=2$ to $L=8$ the length is multiplied by $4$. By what factor is the period
        multiplied? Check the same thing from $L=8$ to $L=32$. What does the exponent $\frac12$
        hidden in the square root do to the growth? \ansline{$\pi\div\frac{\pi}{2}=2$, so the period
        only \emph{doubles}. From $L=8$ to $L=32$ the length again quadruples and
        $2\pi\div\pi=2$ --- doubling again. An exponent of $\frac12$ damps the growth: the output
        grows like the \emph{square root} of the input, so it takes a $4\times$ change in $L$ to buy
        a $2\times$ change in $T$. (Same story as Lesson 5.1's $\frac34$ exponent, and the same
        flattening shape read off the square root graph in 5.0.)}
  \item A clockmaker wants a pendulum that swings exactly \emph{twice} as slowly as her current
        one. She reasons, ``I will just double the length.'' Why is she wrong, and what should she
        do instead? \ansline{Doubling $L$ multiplies $T$ by $\sqrt2\approx1.41$, not by $2$ --- only
        about a $41\%$ slowdown. To double the period she must \textbf{quadruple} the length,
        because $\sqrt4=2$.}
  \item Show that $\sqrt{\dfrac{L}{32}}$ can also be written $\dfrac{\sqrt{2L}}{8}$, so that
        $T(L)=\dfrac{\pi\sqrt{2L}}{4}$. \emph{Hint:} multiply the top and bottom \emph{inside} the
        radical by $2$ and look at the new denominator. Verify your formula at $L=18$.
        \ansline{$\sqrt{\dfrac{L}{32}}=\sqrt{\dfrac{2L}{64}}=\dfrac{\sqrt{2L}}{\sqrt{64}}
        =\dfrac{\sqrt{2L}}{8}$ --- doubling top and bottom turned the denominator into a perfect
        square. Then $T=2\pi\cdot\dfrac{\sqrt{2L}}{8}=\dfrac{\pi\sqrt{2L}}{4}$.\\[3pt]
        Check at $L=18$: $\dfrac{\pi\sqrt{36}}{4}=\dfrac{6\pi}{4}=\dfrac{3\pi}{2}$, matching the
        table.}
\end{enumerate}
\end{extensionbox}

\begin{spiralbox}
\textbf{Coming next --- Lesson 5.3: Multiplying \& Dividing Radicals; Rationalizing.} Today you ran
the product property in one direction, pulling perfect powers \emph{out}. Next you run it the other
way --- multiplying radicals together and distributing across sums, including FOIL on things like
$\left(3+\sqrt5\right)\left(2-\sqrt5\right)$. Then you finally clear simplest-form condition 3, the
one we skipped: a radical stuck in a \emph{denominator}. The tool for a binomial denominator is the
\textbf{conjugate} --- the same trick that turned $a+bi$ into a real number back in Lesson 2.4.
\end{spiralbox}

\begin{teachernote}
Items 1--4 are the routine load, about twenty minutes. Spot-check 1(f) and 2(e): a \emph{negative}
radicand under an odd index is legal, and students who internalized 5.0's ``even index refuses
negatives'' sometimes over-apply it and mark these undefined. Item 3(e) and 3(h) are the two that
require a student to \emph{finish simplifying and then still say ``cannot combine''} --- the point is
that ``cannot combine'' has to be earned, not guessed; 3(h) also quietly rewards anyone who notices
$\sqrt[3]{8}=2$ is not a radical at all.

\textbf{Item 5(b) is the best diagnostic on the page.} A student who says $\sqrt[3]{9}$ is not
simplified because ``$9$ is a perfect square'' has not internalized that the \emph{index} sets the
group size --- the single most common structural misunderstanding of this lesson, and the one that
breaks 5.3's multiplication work first.

\textbf{Item 6 is the conceptual spine} and worth grading for the explanation, not the answer.
``Divide the exponent by the index'' is a rule students will follow correctly and understand
poorly; naming it as improper-fraction-to-mixed-number ties it to Lesson 5.1 and to Warm-Up item 3,
and makes it reconstructible in April. Accept any wording that connects $\frac92=4\frac12$ to
``$x^{4}$ out, one square root left in.''

The \textbf{Extension} is the pendulum, and its shape deliberately matches the Activity's Kepler
problem in 5.1 and the free-fall model in 5.0: a radical model, a table of values chosen to come
out exactly, a growth-rate interpretation, and a structural rewrite. Part (c)--(d) is the payoff ---
the square root's $\frac12$ exponent is why $4\times$ the length buys only $2\times$ the period, and
why the clockmaker's instinct fails. Part (e) is a legitimate preview of 5.3: making the
denominator a perfect square by multiplying inside the radical is rationalizing in its simplest
form, done before the word exists.
\end{teachernote}

\end{document}
